\chapter{Методика}
\label{ch:methods}

%начинать со слов для того, чтобы определить коэффициент Джоуля-Томсона, нужно... из литературы известно... не нужно переписывать учебник
Для того, чтобы определить коэффициент Джоуля-Томсона, нужно исследовать изменение температуры газа при его адиабатическом дросселировании.

Из литературы известно, что коэффициент Джоуля–Томсона определяется, как:

\begin{equation}
\mu = \frac{\Delta T}{\Delta P},
\end{equation}


Для идеального газа энтальпия — однозначная функция температуры. В связи с постоянством энтальпии для идеального газа $\Delta T = 0$, то есть $\mu = 0$. Однако, для реального газа это уже далеко не так, поэтому можно рассмотреть модель газа Ван-дер-Ваальса. Для него получим выражение (вывод см. в Приложении А.):

\begin{equation}
\mu = \frac{\Delta T}{\Delta P} = \frac{\frac{2a}{RT} - b}{C_P}.
\end{equation}

Из этого выражения для $\mu$ сразу видно, что изменение температуры газа Ван-дер-Ваальса при необратимом расширении определяется поправками $a$ и $b$, которые можно найти, зная зависимость $\mu$ от $1/T$. А коэффициент Джоуля-Томсона можно определить по зависимости $\Delta T$ от $\Delta P$.
